% ═══════════════════════════════════════════════════════════════
% ARBEITSBLATT: Reaktion von Metallen mit Sauerstoff
% Thema 2.4 | Chemie Klasse 8 | Kompakt 2 Seiten
% ═══════════════════════════════════════════════════════════════

\documentclass[10pt, a4paper]{article}

\usepackage[utf8]{inputenc}
\usepackage[T1]{fontenc}
\usepackage[ngerman]{babel}

\usepackage{helvet}
\renewcommand{\familydefault}{\sfdefault}

\usepackage{microtype}
\usepackage[margin=1.3cm, top=1.8cm, bottom=1.2cm]{geometry}
\usepackage{enumitem}
\usepackage{tabularx}
\usepackage{booktabs}
\usepackage{array}
\usepackage{multicol}
\usepackage{wrapfig}
\usepackage{tikz}
\usetikzlibrary{calc}

\usepackage{amsmath, amssymb}
\usepackage{siunitx}
\sisetup{locale = DE, per-mode = symbol}

\usepackage{fancyhdr}
\usepackage{lastpage}

\usepackage{xcolor}
\usepackage{tcolorbox}
\tcbuselibrary{skins}

\usepackage{qrcode}

% ═══════════════════════════════════════════════════════════════
% FARBEN
% ═══════════════════════════════════════════════════════════════

\definecolor{chemie}{HTML}{2a7a4b}
\definecolor{hellgrau}{HTML}{F5F5F5}
\definecolor{mittelgrau}{HTML}{888888}
\definecolor{dunkelgrau}{HTML}{444444}

\colorlet{fachfarbe}{chemie}

% ═══════════════════════════════════════════════════════════════
% VARIABLEN
% ═══════════════════════════════════════════════════════════════

\newcommand{\fach}{Chemie}
\newcommand{\thema}{Reaktion von Metallen mit Sauerstoff}
\newcommand{\klasse}{8}
\newcommand{\maxpunkte}{40}

% ═══════════════════════════════════════════════════════════════
% KOPF- UND FUSSZEILE
% ═══════════════════════════════════════════════════════════════

\pagestyle{fancy}
\fancyhf{}
\renewcommand{\headrulewidth}{0.4pt}
\renewcommand{\footrulewidth}{0pt}
\lhead{\small \fach{} Klasse \klasse{} | \thema}
\rhead{\small Name: \underline{\hspace{3cm}}}
\cfoot{\small\textcolor{mittelgrau}{Seite \thepage{} von \pageref{LastPage}}}

% ═══════════════════════════════════════════════════════════════
% BOXEN
% ═══════════════════════════════════════════════════════════════

\newtcolorbox{merkkasten}{
    colback=hellgrau,
    colframe=hellgrau,
    boxrule=0pt,
    arc=0pt,
    left=6pt, right=6pt, top=5pt, bottom=4pt,
    borderline north={2pt}{0pt}{fachfarbe},
    before skip=4pt, after skip=4pt
}

% ═══════════════════════════════════════════════════════════════
% BEFEHLE
% ═══════════════════════════════════════════════════════════════

\newcommand{\luecke}[1]{\underline{\hspace{#1}}}
\newcommand{\punktebox}[1]{\hfill\small\textcolor{mittelgrau}{[\_\_/#1]}}

\newcommand{\aufgabe}[2]{%
    \vspace{6pt}%
    \noindent\textbf{\textcolor{fachfarbe}{Aufgabe #1}} #2%
    \vspace{3pt}%
    \nopagebreak[4]%
}

\newcommand{\operator}[1]{\textbf{\textcolor{fachfarbe}{#1}}}

\newlist{teil}{enumerate}{2}
\setlist[teil,1]{label=\alph*), leftmargin=1.6em, itemsep=2pt, parsep=0pt, topsep=2pt}

\newcommand{\antwortzeilen}[1]{%
    \par\vspace{3pt}%
    \foreach \n in {1,...,#1}{%
        \noindent\rule{\linewidth}{0.3pt}\vspace{5pt}%
    }%
}

\setlength{\parindent}{0pt}
\setlength{\parskip}{2pt}

% ═══════════════════════════════════════════════════════════════
% DOKUMENT
% ═══════════════════════════════════════════════════════════════

\begin{document}

% ═══════════════════════════════════════════════════════════════
% HEADER
% ═══════════════════════════════════════════════════════════════

\begin{center}
    {\large\bfseries\textcolor{fachfarbe}{\thema}}\\[0.1em]
    {\small Arbeitsblatt | \fach{} Klasse \klasse{} | Thema 2.4}
\end{center}
\vspace{-0.4em}
\hrule
\vspace{0.4em}

% ═══════════════════════════════════════════════════════════════
% MERKSÄTZE
% ═══════════════════════════════════════════════════════════════

\begin{merkkasten}
\begin{minipage}[t]{0.85\textwidth}
\textbf{Merksatz 1 – Oxidation:}\\[0.2em]
\luecke{13cm}\\[0.6em]
\luecke{13cm}

\vspace{0.5em}

\textbf{Merksatz 2 – Oxid:}\\[0.2em]
Ein \luecke{2cm} ist eine chemische \luecke{2.5cm} aus einem Element und \luecke{2.5cm}.
\end{minipage}%
\hfill
\begin{minipage}[t]{0.12\textwidth}
\raggedleft
\qrcode[height=1.2cm]{https://devbox42.github.io/sciencesim/chemie/kl08/01-metalloxide/LP-01-metalloxide.html}\\
{\tiny Lernpfad}
\end{minipage}
\end{merkkasten}

% ═══════════════════════════════════════════════════════════════
% AUFGABEN
% ═══════════════════════════════════════════════════════════════

\aufgabe{1}{Begriffe zuordnen} \punktebox{4}

\operator{Ordne} die Begriffe den Definitionen zu.

\small
\begin{tabular}{|p{3cm}|p{10cm}|}
\hline
\textbf{Begriff} & \textbf{Definition} \\
\hline
& Reaktion mit Sauerstoff \\
\hline
& Stoff aus einer Atomsorte \\
\hline
& Stoff aus verschiedenen Atomen \\
\hline
& Verbindung mit Sauerstoff \\
\hline
\end{tabular}

{\scriptsize\textit{Begriffe: Element, Verbindung, Oxid, Oxidation}}

\vspace{6pt}

\aufgabe{2}{Wortgleichungen} \punktebox{6}

\operator{Ergänze} die Wortgleichungen.

\small
\begin{teil}
    \item Magnesium + Sauerstoff → \luecke{4cm}
    \item Eisen + Sauerstoff → \luecke{4cm}
    \item \luecke{3cm} + Sauerstoff → Kupferoxid
    \item Aluminium + Sauerstoff → \luecke{4cm}
    \item Zink + \luecke{3cm} → Zinkoxid
    \item \luecke{3cm} + \luecke{3cm} → Calciumoxid
\end{teil}

\vspace{-4pt}

% ═══════════════════════════════════════════════════════════════
% AUFGABE 3 MIT WERTIGKEITEN
% ═══════════════════════════════════════════════════════════════

\aufgabe{3}{Formelgleichungen aufstellen} \punktebox{6}

\begin{minipage}[t]{0.75\textwidth}
\operator{Schreibe} zu den Wortgleichungen die Formelgleichungen.

\begin{teil}
    \item Magnesium + Sauerstoff → Magnesiumoxid

    \luecke{2.5cm} + \luecke{1.5cm} → \luecke{2.5cm}

    \item Aluminium + Sauerstoff → Aluminiumoxid

    \luecke{2.5cm} + \luecke{1.5cm} → \luecke{2.5cm}

    \item Eisen + Sauerstoff → Eisenoxid (Fe$_2$O$_3$)

    \luecke{2.5cm} + \luecke{1.5cm} → \luecke{2.5cm}
\end{teil}
\end{minipage}%
\hfill
\begin{minipage}[t]{0.22\textwidth}
\raggedleft
\small
\textbf{Wertigkeiten:}\\[0.3em]
\begin{tabular}{|c|c|}
\hline
Mg & II \\
\hline
Al & III \\
\hline
Fe & II/III \\
\hline
Cu & I/II \\
\hline
Zn & II \\
\hline
Ca & II \\
\hline
O & II \\
\hline
\end{tabular}
\end{minipage}

\vspace{6pt}

% ═══════════════════════════════════════════════════════════════
% AUFGABE 4: BEOBACHTUNGEN + FLAMMENFARBEN
% ═══════════════════════════════════════════════════════════════

\aufgabe{4}{Beobachtung und Protokoll} \punktebox{6}

\textbf{a)} Verbrennung von Magnesium:

\small
\begin{tabular}{|p{3cm}|p{10cm}|}
\hline
\textbf{Vor Reaktion} & \\[0.8em]
\hline
\textbf{Während} & \\[0.8em]
\hline
\textbf{Nach Reaktion} & \\[0.8em]
\hline
\textbf{Energie} & \\[0.8em]
\hline
\end{tabular}

\vspace{6pt}

\textbf{b)} Flammenfarben (Simulation): \hfill\qrcode[height=0.9cm]{https://devbox42.github.io/sciencesim/chemie/kl08/01-metalloxide/SIM-01-oxidation.html}

\small
\begin{tabular}{|p{3cm}|p{4cm}|p{4cm}|}
\hline
\textbf{Metall} & \textbf{Flammenfarbe} & \textbf{Oxidfarbe} \\
\hline
Magnesium & & \\[0.6em]
\hline
Eisen & & \\[0.6em]
\hline
Aluminium & & \\[0.6em]
\hline
Kupfer & & \\[0.6em]
\hline
\end{tabular}

\newpage

% ═══════════════════════════════════════════════════════════════
% AUFGABE 5 + 6
% ═══════════════════════════════════════════════════════════════

\aufgabe{5}{Massenerhaltung} \punktebox{8}

\begin{teil}
    \item \operator{Erkläre} das Gesetz von der Erhaltung der Masse.

    \antwortzeilen{2}

    \textbf{Gesetz:} \luecke{10cm}

    \item \operator{Berechne:} 12\,g Magnesium reagieren mit 8\,g Sauerstoff.

    Masse des Magnesiumoxids: \luecke{4cm}
\end{teil}

\vspace{6pt}

\aufgabe{6}{Korrosion im Alltag} \punktebox{4}

\begin{teil}
    \item \operator{Nenne} zwei Beispiele für Korrosion.

    \antwortzeilen{2}

    \item \operator{Erkläre}, warum Aluminium nicht rostet.

    \antwortzeilen{2}
\end{teil}

\vspace{4pt}

% ═══════════════════════════════════════════════════════════════
% AUFGABE 7: VERGLEICH
% ═══════════════════════════════════════════════════════════════

\aufgabe{7}{Vergleich: Magnesium und Aluminium} \punktebox{6}

\operator{Vergleiche} die Oxidation von Magnesium und Aluminium mithilfe der Simulation.

\begin{center}
\small
\begin{tabular}{|p{4cm}|p{4.5cm}|p{4.5cm}|}
\hline
\textbf{Kriterium} & \textbf{Magnesium} & \textbf{Aluminium} \\
\hline
Oxidschicht-Struktur & & \\[0.9em]
\hline
Reaktion stoppt? & & \\[0.9em]
\hline
Fachbegriff & & \\[0.9em]
\hline
\end{tabular}
\end{center}

\operator{Erkläre} den Unterschied in 1--2 Sätzen:

\antwortzeilen{2}

\vspace{6pt}

% ═══════════════════════════════════════════════════════════════
% ZUSATZAUFGABE (unbewertet)
% ═══════════════════════════════════════════════════════════════

\noindent\textcolor{mittelgrau}{\rule{\linewidth}{0.4pt}}

\vspace{4pt}
\noindent{\small\textcolor{mittelgrau}{\textbf{Zusatzaufgabe} (für Schnelle, unbewertet)}}

\vspace{4pt}
\noindent\small\operator{Stelle} die Oxidformeln auf. Nutze die Kreuzregel und die gegebenen Wertigkeiten.

\vspace{4pt}
\small
\begin{tabular}{|p{4cm}|p{2cm}|p{3cm}|}
\hline
\textbf{Name} & \textbf{Wertigkeit} & \textbf{Formel} \\
\hline
Calciumoxid & Ca(II) & \\[0.5em]
\hline
Zinkoxid & Zn(II) & \\[0.5em]
\hline
Natriumoxid & Na(I) & \\[0.5em]
\hline
Bariumoxid & Ba(II) & \\[0.5em]
\hline
\end{tabular}

% ═══════════════════════════════════════════════════════════════
% PUNKTE
% ═══════════════════════════════════════════════════════════════

\vfill
\hrule
\vspace{0.3em}

\hfill\textbf{Punkte:} \luecke{1cm} / \maxpunkte

\end{document}
